% =====================================================================
%  Fiche de lecture (Phase 0) — Projet 1
%  Solveur de Nash m x n et duopoles (Cournot, Bertrand)
%  Groupe 1 — Théorie des Jeux et Systèmes Multi-Agents — ENSA Safi
%  Compilation : pdflatex fiche_lecture.tex   (2 pages maximum)
% =====================================================================
\documentclass[11pt,a4paper]{article}

\usepackage[utf8]{inputenc}
\usepackage[T1]{fontenc}
\usepackage[french]{babel}
\usepackage{amsmath,amssymb}
\usepackage{geometry}
\usepackage{enumitem}
\usepackage{booktabs}
\usepackage{xcolor}
\usepackage{titlesec}

\geometry{margin=1.9cm}
\titlespacing*{\section}{0pt}{6pt}{3pt}
\setlist[itemize]{itemsep=1pt,topsep=2pt,parsep=0pt}

\newcommand{\R}{\mathbb{R}}

\title{\vspace{-1.2cm}\textbf{Fiche de lecture (Phase 0) — Projet 1}\\[2pt]
\large Solveur de Nash $m\times n$ et duopoles (Cournot, Bertrand)}
\author{Groupe 1 : ACHTOUK Fatine, ADIL Youness, BELHIMER Ikram, BOULHAIL Hiba\\
\small Module : Théorie des Jeux et Systèmes Multi-Agents --- ENSA Safi --- Pr.\ A.\ JRAIFI}
\date{Année universitaire 2026--2027 --- Semaine 1}

\begin{document}
\maketitle
\vspace{-0.6cm}

\section*{1. Formules clés}
\noindent\textbf{Conventions.} $A\in\R^{m\times n}$ : gains du joueur 1 (lignes) ;
$B\in\R^{m\times n}$ : gains du joueur 2 (colonnes) ; indices à partir de $0$.
$x\in\Delta_m$ (joueur 1), $y\in\Delta_n$ (joueur 2).

\smallskip
\noindent\textbf{Équilibre de Nash (mixte).} $(x,y)$ est un équilibre de Nash si et seulement
si aucune \emph{déviation pure} n'est profitable :
\begin{equation}
(Ay)_i \le x^{\top} A y \quad \forall i \in \{1,\dots,m\},
\qquad
(x^{\top}B)_j \le x^{\top} B y \quad \forall j \in \{1,\dots,n\}.
\end{equation}
\emph{Avec mes mots :} au lieu de tester toutes les stratégies mixtes (infinité), il suffit de
vérifier que chaque joueur ne gagne pas plus en jouant une stratégie \emph{pure} plutôt que sa
stratégie mixte. C'est cette condition, simple à tester, qu'on utilise pour valider chaque
solution.

\smallskip
\noindent\textbf{Indifférence sur les supports.} Soient $I=\mathrm{supp}(x)$ et
$J=\mathrm{supp}(y)$ avec $|I|=|J|=k$. Sur le support, le joueur 1 est indifférent entre
toutes ses stratégies de $I$, et de même pour le joueur 2 sur $J$ :
\begin{equation}
(Ay)_i = v_1 \ \ (i\in I),
\qquad
(x^{\top}B)_j = v_2 \ \ (j\in J),
\qquad
\sum_{i\in I} x_i = 1,\quad \sum_{j\in J} y_j = 1 .
\end{equation}
\emph{Avec mes mots :} un joueur qui mélange plusieurs stratégies doit être exactement
indifférent entre elles (sinon il ne jouerait que la meilleure). Cela donne un petit système
linéaire à résoudre pour chaque couple de supports.

\smallskip
\noindent\textbf{Cournot} ($P(Q)=a-bQ$, coût marginal constant $c$, $a>c$) :
\begin{equation}
\text{meilleure réponse : } q_i=\frac{a-c-b\sum_{j\ne i} q_j}{2b},
\qquad
q^{*}=\frac{a-c}{(n+1)b},
\qquad
\pi^{*}=\frac{(a-c)^2}{b\,(n+1)^2}.
\end{equation}
Monopole : $Q=\dfrac{a-c}{2b}$, $\pi=\dfrac{(a-c)^2}{4b}$.
Duopole ($n=2$) : $q^{*}=\dfrac{a-c}{3b}$, $\pi^{*}=\dfrac{(a-c)^2}{9b}$.
Cartel (deux firmes) : $q=\dfrac{a-c}{4b}$, $\pi=\dfrac{(a-c)^2}{8b}$.

\smallskip
\noindent\textbf{Bertrand} (biens homogènes, coûts marginaux égaux $c$) : l'unique équilibre
est $p_1=p_2=c$ et le profit est nul (paradoxe de Bertrand), alors que le monopole fixerait
$p^{M}=\dfrac{a+c}{2}$.

\section*{2. Lien avec le code du cours}
\begin{itemize}
  \item \texttt{nash\_mixte\_2x2} (cours) $\rightarrow$ \textbf{à généraliser} en
        \texttt{enumeration\_supports(A, B, tol=1e-9)} : la version du cours est limitée à
        $2\times2$, la nôtre balaie tous les supports $k=1,\dots,\min(m,n)$.
  \item Nouvelle fonction de contrôle \texttt{est\_nash(A, B, x, y, tol=1e-9)} : teste (1)
        pour tout profil candidat (elle sert aussi de garde-fou dans les tests).
  \item Nouvelle fonction \texttt{cournot(a, b, c, n=2)} (formules fermées) et
        \texttt{cournot\_grille(a, b, c, qmax)} qui construit les matrices de gains et
        réutilise \texttt{equilibres\_nash\_purs} du cours.
  \item \textbf{Outils agent à ajouter} dans \texttt{agent\_tj\_sma.py} :
        \texttt{resoudre\_bimatrice(A, B)} et \texttt{cournot(a, b, c, n\_firmes)}
        (plafonds sur $m$, $n$ et le nombre de firmes ; sortie JSON ; types annotés).
\end{itemize}

\section*{3. Algorithme : démarche de calcul en Python}
\begin{enumerate}[label=\arabic*.]
  \item Pour $k=1$ à $\min(m,n)$ :
  \item énumérer chaque support $I$ de lignes ($|I|=k$) et chaque support $J$ de colonnes
        ($|J|=k$) ;
  \item résoudre le système (2) sur $(I,J)$ : inconnues $(y_j)_{j\in J}$ et $v_1$, puis
        $(x_i)_{i\in I}$ et $v_2$ (deux petits systèmes linéaires indépendants,
        \texttt{numpy.linalg.lstsq}) ;
  \item rejeter si une composante est négative (hors du simplexe) ou si le système est
        incompatible (jeu dégénéré) ;
  \item tester la non-déviation (1) via \texttt{est\_nash} ; conserver l'équilibre sinon.
\end{enumerate}
\emph{Complexité :} le nombre de couples de supports vaut
$\sum_{k} \binom{m}{k}\binom{n}{k}$, soit une croissance exponentielle en la taille --- d'où
l'expérience « nombre d'équilibres et temps de calcul en fonction de la taille ». Pour Cournot,
on compare $q^{*}$ (formule fermée) aux équilibres purs obtenus sur une grille de quantités.

\section*{4. Hypothèses du modèle et limites}
\begin{itemize}
  \item Jeux \emph{statiques}, finis, à \emph{information complète} ; rationalité commune.
  \item \textbf{Jeux non dégénérés} pour l'énumération des supports : dans un jeu dégénéré, des
        équilibres à supports de tailles différentes existent et le système d'indifférence
        devient sous-déterminé (à signaler, pas à conclure « pas d'équilibre »).
  \item \textbf{Cournot} : demande linéaire, coûts marginaux constants, quantités choisies
        simultanément, produit homogène. \textbf{Bertrand} : biens homogènes, coûts égaux,
        prix simultanés.
  \item Limite numérique : toutes les comparaisons se font avec une \emph{tolérance}
        (\texttt{tol=1e-9}) --- jamais d'égalité de flottants exacte.
\end{itemize}

\section*{5. Données et paramètres de référence}
\begin{itemize}
  \item Jeux du cours : Jeu 1, Jeu 2, dilemme du prisonnier, bataille des sexes
        (mixte $p=\tfrac23$, $q=\tfrac13$), pierre-feuille-ciseaux ($3\times3$, équilibre
        $(\tfrac13,\tfrac13,\tfrac13)$) ; jeux aléatoires $m\times n$ à graine fixée.
  \item Cournot : $a=100$, $b=1$, $c=10$ $\Rightarrow$ $q^{*}=30$ par firme et profit $900$ ;
        cartel : $22{,}5$ par firme et profit $1012{,}5$.
  \item Reproductibilité : graines fixées pour toutes les expériences
        (\texttt{np.random.default\_rng(seed)}).
\end{itemize}

\end{document}
